\documentclass[11pt]{article}

\usepackage[margin=1in]{geometry}
\usepackage{amsmath, amssymb, amsthm, amsfonts}
\usepackage{mathtools}
\usepackage{bm}
\usepackage{bbm}
\usepackage{hyperref}
\usepackage{enumitem}
\usepackage{physics}

\hypersetup{
  colorlinks=true,
  linkcolor=blue,
  citecolor=blue,
  urlcolor=blue
}

\title{Prime-Indexed Tensor Dynamics for \\ Quantum Agentic Governance Intelligence:\\
An Operator-Theoretic Framework with Testable Constraints}

\author{Citizen Gardens \\ The Foundation of Multiplicity}

\date{April 25, 2026}

\begin{document}

\maketitle

\begin{abstract}
Prime-Indexed Tensor Networks (PITNs) are proposed as a tensor-network substrate for Quantum Agentic Governance Intelligence (QAGI). Rather than treating intelligence as a static function approximator, PITNs model cognition as an evolving tensor-network state on a constrained Hilbert-space manifold, driven by prime-indexed operator sequences and equipped with intrinsic execution attestation. The framework replaces earlier ``guaranteed convergence'' language with operator-controlled dynamics under explicit norm and projection assumptions, making the claims mathematically defensible and experimentally verifiable. We formalize the state space, the prime-indexed operator grammar, a spectral anchor and ecological memory term, and a Prime-Weighted Execution Hashing (PWEH) integrity functional. We then state one provable proposition on norm stability and two testable conjectures on sample efficiency and trace-based integrity. Connections to time-dependent variational principle (TDVP) methods for tensor networks and to post-quantum cryptographic hash functions are highlighted to ground the architecture in existing physics and cryptography literature.\footnote{For TDVP and tensor-network evolution on constrained manifolds, see e.g.\ TDVP implementations for MPS and TTN.[web:51][web:76][web:79] For tensor-network based quantum machine learning, see.[web:80] For post-quantum hash and symmetric-primitive assumptions, see.[web:86][web:108]} 
\end{abstract}
\newpage
\tableofcontents
\newpage
\section{Executive Summary}

This report develops Prime-Indexed Tensor Networks (PITNs) from a conceptual model into a mathematically well-posed substrate for Quantum Artificial General Intelligence (QAGI). The central move is to treat PITNs as \emph{evolving tensor-network states} on a constrained variational manifold, governed by a discrete alphabet of prime-indexed operators rather than by fully associative layer compositions.

\medskip

The key elements are:
\begin{itemize}[leftmargin=*]
  \item A composite Hilbert space \( \mathcal{H} = \bigotimes_{k=1}^K \mathcal{H}_k \) and a variational tensor manifold \( \mathcal{M} \subset \mathcal{H} \) with bounded bond dimension, analogous to TDVP-style tensor-network manifolds.[web:51][web:79][web:80]
  \item A time-dependent active set of primes \( P_N(t) \subset \mathbb{P} \) and a family of bounded prime operators \( A_p \in \mathcal{B}(\mathcal{H}) \), whose ordered compositions define a non-associative, path-dependent grammar over network construction.
  \item A spectral anchor operator \( \Lambda_m \) (projection or contraction onto \( \mathcal{M} \)) and an ecological memory term \( Q_{ent} \) that couples the current state to a finite history window via entanglement-aware observables.
  \item A Prime-Weighted Execution Hash (PWEH) functional \( S_{\text{integrity}} \) that accumulates a post-quantum cryptographic hash over the ordered prime trace and local state norms, providing intrinsic execution attestation.[web:86][web:108]
\end{itemize}

We deliberately distinguish:
\begin{enumerate}[leftmargin=*]
  \item \textbf{Provable properties} under bounded-operator and projection assumptions, including norm-bounded evolution and local stability of fixed points.
  \item \textbf{Testable conjectures} regarding multiplicative vs.\ positional indexing efficiency and the robustness of trace-based integrity against adversarial perturbations.
\end{enumerate}

This separation reframes the PITN/QAGI narrative as a rigorous machine-learning architecture with a clear experimental agenda, rather than as a set of unqualified AGI guarantees.

\section{Background and Motivation}

Tensor networks provide a structured way to represent high-dimensional states in quantum many-body physics and, increasingly, in machine learning.[web:80][web:100] Algorithms such as DMRG, TEBD, and the Time-Dependent Variational Principle (TDVP) evolve states on manifolds of Matrix Product States (MPS), Projected Entangled Pair States (PEPS), or Tree Tensor Networks (TTN) while controlling entanglement growth and computational complexity.[web:51][web:76][web:79][web:104]

In parallel, post-quantum cryptography (PQC) has developed hash-based, lattice-based, and other schemes designed to remain secure against quantum adversaries.[web:86][web:108][web:101] Hash-based constructions, in particular, provide conceptually simple integrity mechanisms that can serve as building blocks for signatures and attestation.

PITNs aim to unify these threads in a QAGI context:
\begin{itemize}[leftmargin=*]
  \item \emph{From tensor networks}, we borrow the idea of evolving a high-dimensional state on a constrained variational manifold with explicit control of bond dimension and entanglement.
  \item \emph{From prime number theory}, we borrow a multiplicative indexing scheme that enables unique factorization of interaction pathways and makes the order of prime-channel application a first-class object of learning.
  \item \emph{From PQC}, we borrow execution attestation ideas, using prime-weighted operator traces as inputs to post-quantum hash functions.
\end{itemize}

The resulting architecture is both structurally novel (prime-indexed non-associative grammar) and aligned with established operator and tensor-network methods.

\section{Formal Specification: Prime-Indexed Tensor Dynamics}

\subsection{State space and variational manifold}

Let \( \mathcal{H} = \bigotimes_{k=1}^K \mathcal{H}_k \) be a composite Hilbert space. We consider a variational tensor manifold \( \mathcal{M} \subset \mathcal{H} \) defined, for example, by tensor-network states (MPS, TTN, etc.) with a maximum bond dimension \( \chi \).[web:80][web:100] 

\begin{description}[leftmargin=2em]
\item[Assumption 1 (Smooth variational manifold).] 
\( \mathcal{M} \) is a smooth submanifold of \( \mathcal{H} \), endowed with the inner product inherited from \( \mathcal{H} \). The tangent space \( T_T\mathcal{M} \) is well-defined for all \( T \in \mathcal{M} \).

\item[State variable.] 
The PITN state at discrete time \( t \in \mathbb{N} \) is a tensor-network state \( T(t) \in \mathcal{M} \subset \mathcal{H} \).

\end{description}

This mirrors TDVP-like formulations where time evolution is projected back onto the manifold, typically via a tangent-space projector that ensures the state stays within the chosen tensor-network class.[web:51][web:76][web:79]

\subsection{Admissible prime operators and ordered grammar}

Let \( \mathbb{P} \) denote the set of prime numbers, and let \( P_N(t) \subset \mathbb{P} \) be a finite, time-dependent active prime set.

\begin{description}[leftmargin=2em]
\item[Definition 1 (Prime operators).]
For each \( p \in P_N(t) \), define an admissible prime operator
\[
A_p \in \mathcal{B}(\mathcal{H}),
\]
where \( \mathcal{B}(\mathcal{H}) \) is the space of bounded linear operators on \( \mathcal{H} \), equipped with the operator norm
\[
\|A\|_{\mathrm{op}} = \sup_{\|x\|=1} \|Ax\|.
\][web:82][web:87]

\item[Prime words and composition.] 
Given a sequence of prime indices (a \emph{prime word})
\[
w(t) = (p_{i_1}, p_{i_2}, \dots, p_{i_t}),
\]
define the associated composite transformation
\[
\mathcal{A}_{w(t)} = A_{p_{i_t}} A_{p_{i_{t-1}}} \cdots A_{p_{i_1}}.
\]

\end{description}

While multiplication in \( \mathcal{B}(\mathcal{H}) \) is associative, the \emph{effective} PITN dynamics are not, due to interleaved projections and coupling to ecological memory.

\begin{description}[leftmargin=2em]
\item[Non-associativity and path dependence.] 
Let \( \Lambda \) denote the spectral anchor (defined below) and \( Q_{ent} \) the ecological memory term. Define the one-step update as
\[
T(t+1) = \Lambda\big(A_{p_{i_t}} T(t) + Q_{ent}(t)\big).
\]
Because each application of \( A_{p_{i_t}} \) is followed by \( \Lambda \) and \( Q_{ent} \) depends on the trajectory history, the effective transformations
\[
(\Lambda A_{p_2} \Lambda A_{p_1})T
\quad\text{and}\quad
(\Lambda A_{p_1} \Lambda A_{p_2})T
\]
need not coincide. The state \( T(t) \) is thus inherently path-dependent: it encodes not only the data seen but also the ordered sequence of prime-indexed operations used to process it, analogously to path-dependent dynamics in non-ergodic reinforcement learning.[web:88]
\end{description}

\subsection{Spectral anchor and ecological memory}

\subsubsection{Spectral anchor \( \Lambda_m \)}

The spectral anchor \( \Lambda_m \) maps the intermediate result of prime operations back onto the variational manifold and controls the growth of the state norm.

\begin{description}[leftmargin=2em]
\item[Definition 2 (Spectral anchor).] 
For each time step \( t \), let
\[
\Lambda_m : \mathcal{H} \to \mathcal{M}
\]
be an orthogonal projection or, more generally, a contraction operator whose image is contained in \( \mathcal{M} \), satisfying
\[
\|\Lambda_m\|_{\mathrm{op}} \le 1.
\]
In the simplest case, \( \Lambda_m \equiv \Pi_{\mathcal{M}} \) is the orthogonal projection onto \( \mathcal{M} \).

\end{description}

This definition is in line with TDVP approaches where the dynamics are projected onto the tangent space or manifold to maintain a representation within the chosen tensor-network class, with controlled norm and entanglement growth.[web:51][web:76][web:79]

\subsubsection{Ecological memory \( Q_{ent} \)}

The ecological memory term \( Q_{ent} \) couples the current state to a history window via entanglement-aware observables.

\begin{description}[leftmargin=2em]
\item[Definition 3 (Ecological memory).]
Let \( X_{t-j} \) denote exogenous inputs or environmental states at past times \( t-j \), and let \( w(t-j) \) denote the prime word up to time \( t-j \). For a history length \( \tau \) and a decay kernel \( \gamma_j \ge 0 \), define
\[
Q_{ent}(t)
=
\sum_{j=1}^{\tau}
\gamma_j\,
\mathcal{E}_{ent}\big(T(t) \otimes A_{w(t-j)}(X_{t-j})\big),
\]
where \( \mathcal{E}_{ent} \) is a non-local functional that evaluates entanglement-aware observables on the composite tensor network.

\end{description}

In practice, \( \mathcal{E}_{ent} \) may be instantiated using measures of entanglement entropy or related observables that quantify non-local correlations across subsystems in a tensor network, as is standard in tensor-network simulations of many-body systems and quantum machine learning.[web:56][web:80][web:104] The delayed dependence on \( X_{t-j} \) and \( w(t-j) \) creates a form of ecological memory and contextual prior injection.

\subsection{Prime-Weighted Execution Hashing (PWEH)}

To provide intrinsic execution attestation, we define a Prime-Weighted Execution Hash that accumulates a post-quantum cryptographic hash over the ordered prime trace and local state norms.

\begin{description}[leftmargin=2em]
\item[Definition 4 (Integrity functional).]
Let \( S_{\text{integrity}}(t) \) denote the integrity state at time \( t \), and let \( \|\cdot\| \) denote a chosen norm on \( \mathcal{H} \). Define
\[
S_{\text{integrity}}(t)
=
Hash_{PQC}\Big(
S_{\text{integrity}}(t-1)
\parallel p_{i_t}
\parallel \|A_{p_{i_t}} T(t)\|
\Big),
\]
where \( Hash_{PQC} \) is a candidate post-quantum cryptographic hash function (e.g., hash-based or lattice-based constructions standardized for PQC) over an appropriate finite field, designed to resist both classical and quantum attacks.[web:86][web:108][web:101]

\end{description}

Here \( \parallel \) denotes concatenation of bitstrings. Because the input depends both on the discrete prime index \( p_{i_t} \) and on the norm of the local update \( A_{p_{i_t}} T(t) \), any significant change in execution order or state magnitude will alter the hash chain, subject to the sensitivity of the chosen encoding and hash function.

\section{Update Dynamics and Stability}

\subsection{Recursive update rule}

Putting the above pieces together, a single PITN update step can be written as
\begin{equation}
  T(t+1)
  =
  \Lambda_m\big(A_{p_{i_t}} T(t) + Q_{ent}(t)\big),
  \label{eq:update}
\end{equation}
where:
\begin{itemize}[leftmargin=*]
  \item \( A_{p_{i_t}} \) is the selected prime operator at time \( t \),
  \item \( Q_{ent}(t) \) is the ecological memory term,
  \item \( \Lambda_m \) projects the result back into \( \mathcal{M} \).
\end{itemize}

The learning process (e.g., in a reinforcement-learning or supervised-learning context) adjusts the parameters of the operators \( A_p \), the decay kernel \( \gamma_j \), the choice of active primes \( P_N(t) \), and possibly the manifold \( \mathcal{M} \) itself (e.g., via bond dimension schedules), all subject to constraints on operator norms and projections.

\subsection{Norm stability under bounded operators}

To avoid catastrophic divergence, we impose bounds on the operator norms and on the ecological memory term.

\begin{description}[leftmargin=2em]
\item[Assumption 2 (Operator and memory bounds).]
\begin{enumerate}[label=(\alph*), leftmargin=2.5em]
  \item For all active \( p \in P_N(t) \), the prime operators are bounded with
  \[
  \|A_p\|_{\mathrm{op}} \le c
  \]
  for some constant \( c > 0 \).[web:82][web:87]
  \item The spectral anchor satisfies \( \|\Lambda_m\|_{\mathrm{op}} \le 1 \).
  \item The ecological memory term satisfies a uniform bound
  \[
  \|Q_{ent}(t)\| \le q_{\max}
  \]
  for all \( t \).
\end{enumerate}
\end{description}

We then have the following basic stability result.

\begin{proposition}[Local norm stability under bounded projection]
\label{prop:norm-stability}
Under Assumption~2, the update rule~\eqref{eq:update} satisfies
\[
\|T(t+1)\| \le c\,\|T(t)\| + q_{\max}.
\]
In particular, if \( c < 1 \), the sequence \( \{\|T(t)\|\}_{t\ge 0} \) is uniformly bounded and converges to a ball of radius \( \frac{q_{\max}}{1-c} \) in norm.
\end{proposition}

\begin{proof}
Using the operator norm inequality and \( \|\Lambda_m\|_{\mathrm{op}} \le 1 \), we have
\[
\|T(t+1)\|
= \|\Lambda_m(A_{p_{i_t}} T(t) + Q_{ent}(t))\|
\le \|A_{p_{i_t}} T(t) + Q_{ent}(t)\|
\le \|A_{p_{i_t}}\|_{\mathrm{op}} \|T(t)\| + \|Q_{ent}(t)\|.
\]
By Assumption~2, this yields
\[
\|T(t+1)\| \le c\,\|T(t)\| + q_{\max}.
\]
Standard iteration of this affine recurrence implies uniform boundedness and convergence to a ball of radius \( \frac{q_{\max}}{1-c} \) when \( c < 1 \).
\end{proof}

This proposition provides a minimal but rigorous guarantee that, under appropriate norm control and bounded ecological memory, PITN dynamics do not diverge arbitrarily in norm. Linearizing around a fixed point \( \bar T \), and assuming local Lipschitz properties, yields a Jacobian with spectral radius controlled by \( c \), implying local fixed-point stability for \( c < 1 \), in line with standard results on bounded linear operators on Hilbert spaces.[web:87]

\section{Conjectures and Experimental Agenda}

\subsection{Conjecture 1: multiplicative vs.\ positional efficiency}

\begin{conjecture}[Multiplicative indexing and sample efficiency]
For hierarchical, highly compositional reinforcement-learning or sequence-modeling tasks, a tensor network governed by prime-indexed multiplicative interactions \( A_p \) and constrained to a manifold \( \mathcal{M} \) with fixed bond dimension \( \chi \) can achieve a target generalization loss \( L_{\text{target}} \) in fewer parameter updates than an architecture with comparable parameter count that uses strictly associative, positional interactions (e.g., a Transformer with positional encodings and no explicit prime grammar).
\end{conjecture}

This conjecture is motivated by:
\begin{itemize}[leftmargin=*]
  \item Results where structured tensor-network ans\"atze efficiently represent certain classes of quantum and classical states that are costly for unstructured models.[web:80][web:104][web:98]
  \item The hypothesis that prime-indexed channels provide a multiplicative basis of interaction modes, enabling reuse of prime patterns across tasks and hierarchical compositions.
\end{itemize}

\paragraph{Validation path.}
\begin{enumerate}[leftmargin=*]
  \item Construct PITN models with a fixed bond dimension \( \chi \) and a controlled active prime set \( P_N(t) \).
  \item Construct baseline models (e.g., Transformers or RNNs) with similar parameter counts and depth.
  \item Evaluate on hierarchical RL tasks or compositional sequence tasks (e.g., algorithmic reasoning, hierarchical control) and compare the number of updates and wall-clock time needed to reach \( L_{\text{target}} \).[web:80][web:95]
\end{enumerate}

\subsection{Conjecture 2: trace-based integrity via PWEH}

\begin{conjecture}[Trace-based integrity under PQC assumptions]
Assume that \( Hash_{PQC} \) satisfies standard post-quantum security assumptions (preimage and second-preimage resistance, collision resistance up to negligible probability) at an appropriate security level.[web:86][web:108][web:101] Further assume that any adversarial perturbation \( \delta T \) to the tensor state or manipulation of the update order \( w(t) \) that significantly alters the predictive output induces a non-negligible change in the encoded hash input. Then, with probability \( 1 - \text{negl} \), such perturbations produce a hash mismatch in \( S_{\text{integrity}}(t) \).
\end{conjecture}

In other words, under both:
\begin{itemize}[leftmargin=*]
  \item cryptographic assumptions on \( Hash_{PQC} \), and
  \item architectural assumptions relating significant output changes to significant changes in the hash input space,
\end{itemize}
PWEH functions as a robust integrity-attestation mechanism for PITN trajectories.

\paragraph{Validation path.}
\begin{enumerate}[leftmargin=*]
  \item Select concrete PQC hash candidates (e.g., hash-based schemes aligned with PQC standards).[web:86][web:108]
  \item Implement PWEH over PITN training and deployment traces.
  \item Conduct adversarial experiments (including model poisoning and manipulation of the prime sequence \( w(t) \)) and empirically measure the probability of undetected perturbations.
  \item Pursue formal cryptanalysis mapping perturbations in \( \mathcal{M} \) to the induced changes in the hash input domain for the chosen scheme.[web:93][web:97][web:101]
\end{enumerate}

\section{Implementation Sketch and Code Snippet}

\subsection{High-level algorithm}

A minimal PITN-based QAGI core loop can be sketched as follows:

\begin{enumerate}[leftmargin=*]
  \item Initialize state \( T(0) \in \mathcal{M} \), integrity state \( S_{\text{integrity}}(0) \), and active prime set \( P_N(0) \).
  \item For each time step \( t = 0,1,\dots \):
  \begin{enumerate}[label=(\alph*), leftmargin=2em]
    \item Observe current input \( X_t \) and ecological context.
    \item Select a prime \( p_{i_t} \in P_N(t) \) (e.g., via a policy \( \pi(p \mid T(t), X_t) \)).
    \item Compute intermediate update \( U(t) = A_{p_{i_t}} T(t) + Q_{ent}(t) \).
    \item Project back onto the manifold: \( T(t+1) = \Lambda_m(U(t)) \).
    \item Update integrity state:
    \[
    S_{\text{integrity}}(t+1)
    =
    Hash_{PQC}\big(S_{\text{integrity}}(t) \parallel p_{i_t} \parallel \|U(t)\|\big).
    \]
    \item Evaluate loss or reward and perform meta-learning over parameters of \( A_p \), \( \Lambda_m \), \( Q_{ent} \), and possibly \( P_N(t) \).
  \end{enumerate}
\end{enumerate}

\subsection{Illustrative pseudo-code (Python style)}

Inspired by TDVP-style implementations for time evolution of tensor networks,[web:62][web:96] an abstract PITN loop might look as follows:

\begin{verbatim}
class PITNCore:
    def __init__(self, manifold, prime_ops, anchor, entanglement_fn,
                 hash_fn, active_primes):
        self.M = manifold            # Variational tensor manifold M
        self.A = prime_ops           # dict: p -> A_p
        self.Lambda = anchor         # spectral anchor projection
        self.E_ent = entanglement_fn # ecological/entanglement functional
        self.hash_fn = hash_fn       # post-quantum hash candidate
        self.P = active_primes       # active prime set P_N(t)
        self.S_integrity = init_hash_state()
        self.state = self.M.random_state()

    def step(self, X_t, history):
        # 1. Select prime according to current policy
        p = select_prime(self.state, X_t, self.P)

        # 2. Build ecological memory term Q_ent(t)
        Q_ent_t = 0
        for j, (T_past, X_past, w_past, gamma_j) in enumerate(history):
            Q_ent_t += gamma_j * self.E_ent(self.state, apply_word(w_past, X_past))

        # 3. Apply prime operator and spectral anchor
        U = self.A[p](self.state) + Q_ent_t
        T_next = self.Lambda(U)   # Project back to M

        # 4. Update integrity hash
        self.S_integrity = self.hash_fn(self.S_integrity, p, norm(U))

        # 5. Commit new state
        self.state = T_next
        return T_next, self.S_integrity
\end{verbatim}

This snippet is schematic but captures the essential ingredients: prime selection, ecological memory construction, manifold projection, and hash-chain update, mirroring the structure of TDVP engines that evolve tensor-network states under constrained dynamics.[web:51][web:96]

\section{Discussion and Outlook}

We have articulated a rigorous operator-theoretic formulation of Prime-Indexed Tensor Networks as a QAGI substrate:
\begin{itemize}[leftmargin=*]
  \item The state lives on a constrained tensor-network manifold \( \mathcal{M} \) with controlled bond dimension and entanglement; evolution is governed by prime-indexed operators and projected through a spectral anchor, in analogy with TDVP methods.[web:51][web:76][web:79][web:80]
  \item The prime grammar is inherently path-dependent and non-associative at the level of effective dynamics, enabling a ``prime move'' view of cognition where intelligence unfolds as sequences of prime-labeled operations on a structured state space.
  \item The ecological memory term \( Q_{ent} \) injects delayed, entanglement-aware context, and the PWEH integrity functional \( S_{\text{integrity}} \) offers an intrinsic, PQC-aligned execution attestation layer.[web:80][web:86][web:108]
\end{itemize}

At the same time, we have explicitly separated provable norm-stability properties from aspirational conjectures. The next steps are clear:
\begin{enumerate}[leftmargin=*]
  \item Implement PITN dynamics on practical tensor-network backends (e.g., MPS/TTN libraries) to test stability and representation power on synthetic and real tasks.[web:62][web:96][web:80]
  \item Benchmark multiplicative vs.\ positional indexing for hierarchical RL and sequence tasks, under equal parameter budgets.
  \item Integrate concrete PQC hash primitives into PWEH and subject them to adversarial testing and formal cryptanalysis.
\end{enumerate}

By grounding PITNs in established operator and tensor-network theory while isolating genuinely novel claims as conjectures, this framework provides a realistic path from theoretical elegance to empirical evidence.

\appendix
\section*{Mathematical Appendix}
\addcontentsline{toc}{section}{Mathematical Appendix}

\section{Operator-Theoretic Preliminaries}

In this appendix we collect and prove the basic operator-norm inequalities and stability bounds used in the main text. Throughout, \( \mathcal{H} \) denotes a complex Hilbert space with norm \( \|\cdot\| \) induced by its inner product, and \( \mathcal{B}(\mathcal{H}) \) denotes the Banach algebra of bounded linear operators on \( \mathcal{H} \) endowed with the operator norm.[web:87][web:119][web:124]

\subsection{Operator norm and submultiplicativity}

\begin{definition}[Operator norm]
Let \( A \in \mathcal{B}(\mathcal{H}) \). The operator norm of \( A \) is defined by
\[
\|A\|_{\mathrm{op}} := \sup_{\|x\| = 1} \|Ax\|.
\]
Equivalently, \( \|A\|_{\mathrm{op}} \) is the least constant \( c \ge 0 \) such that \( \|Ax\| \le c\|x\| \) for all \( x \in \mathcal{H} \).[web:82][web:119]
\end{definition}

\begin{lemma}[Submultiplicativity]\label{lem:submult}
If \( A,B \in \mathcal{B}(\mathcal{H}) \), then
\[
\|BA\|_{\mathrm{op}} \le \|B\|_{\mathrm{op}} \, \|A\|_{\mathrm{op}}.
\]
\end{lemma}

\begin{proof}
For any \( x \in \mathcal{H} \) with \( \|x\| = 1 \), we have
\[
\|BAx\| \le \|B\|_{\mathrm{op}} \, \|Ax\| \le \|B\|_{\mathrm{op}} \, \|A\|_{\mathrm{op}} \, \|x\|
= \|B\|_{\mathrm{op}} \, \|A\|_{\mathrm{op}}.
\]
Taking the supremum over all unit vectors \( x \) yields the result.[web:82][web:87]
\end{proof}

\subsection{Orthogonal projections and contractions}

We recall the basic facts about projections and contractions on Hilbert spaces.[web:87][web:120]

\begin{definition}[Projection]
A linear operator \( P:\mathcal{H} \to \mathcal{H} \) is a projection if \( P^2 = P \). It is an \emph{orthogonal projection} if, in addition, it is self-adjoint, i.e.\ \( P^\ast = P \).[web:87][web:119]
\end{definition}

\begin{theorem}[Projections and direct sum]\label{thm:projection}
If \( P:\mathcal{H}\to\mathcal{H} \) is an orthogonal projection, then
\[
\mathcal{H} = \mathrm{ran}\,P \oplus \ker P,
\]
and both \( \mathrm{ran}\,P \) and \( \ker P \) are closed subspaces.[web:87]
\end{theorem}

\begin{proposition}[Norm of an orthogonal projection]\label{prop:projection-norm}
If \( P \) is a nonzero orthogonal projection on \( \mathcal{H} \), then
\[
\|P\|_{\mathrm{op}} = 1.
\]
\end{proposition}

\begin{proof}
For any \( x \in \mathcal{H} \),
\[
\|Px\|^2 = \langle Px, Px \rangle = \langle x, P^\ast Px \rangle = \langle x, Px \rangle \le \|x\|\,\|Px\|
\]
by Cauchy--Schwarz, so \( \|Px\| \le \|x\| \), and hence \( \|P\|_{\mathrm{op}} \le 1 \).[web:87][web:122]
Since \( P \) is nonzero, there exists \( x \neq 0 \) with \( Px \neq 0 \). For such \( x \),
\[
\|Px\| = \|P(Px)\| \le \|P\|_{\mathrm{op}} \,\|Px\|.
\]
If \( \|P\|_{\mathrm{op}} < 1 \), then this implies \( \|Px\| = 0 \), a contradiction. Thus \( \|P\|_{\mathrm{op}} \ge 1 \). Combining with \( \|P\|_{\mathrm{op}} \le 1 \) yields \( \|P\|_{\mathrm{op}} = 1 \).[web:87]
\end{proof}

\begin{definition}[Contraction]
A bounded operator \( T:\mathcal{H}\to\mathcal{H} \) is called a contraction if \( \|T\|_{\mathrm{op}} \le 1 \).[web:120]
\end{definition}

Every orthogonal projection is a contraction. More generally, any linear operator \( \Lambda_m:\mathcal{H} \to \mathcal{H} \) with \( \|\Lambda_m\|_{\mathrm{op}} \le 1 \) is a contraction and will be used as a \emph{spectral anchor} in the PITN dynamics.

\section{Norm Bounds for PITN Updates}

We now derive the norm bounds used in the main text to establish local stability of the PITN update rule under bounded operators and memory.

\subsection{Setup and assumptions}

Recall the PITN one-step update
\begin{equation}\label{eq:appendix-update}
T(t+1) = \Lambda_m\big(A_{p_{i_t}} T(t) + Q_{ent}(t)\big),
\end{equation}
where:
\begin{itemize}[leftmargin=*]
  \item \( T(t) \in \mathcal{H} \) is the state at time \( t \);
  \item \( A_{p_{i_t}} \in \mathcal{B}(\mathcal{H}) \) is the prime operator at step \( t \);
  \item \( \Lambda_m:\mathcal{H}\to\mathcal{H} \) is the spectral anchor (projection or contraction onto the variational manifold \( \mathcal{M} \subset \mathcal{H} \));
  \item \( Q_{ent}(t) \in \mathcal{H} \) is the ecological memory term.
\end{itemize}

We repeat the assumptions in a form suitable for analysis.

\begin{description}[leftmargin=2em]
\item[Assumption A1 (Bounded prime operators).]
There exists a constant \( c > 0 \) such that for all active primes \( p \in P_N(t) \) and all \( t \),
\[
\|A_p\|_{\mathrm{op}} \le c.
\]

\item[Assumption A2 (Contractive spectral anchor).]
The anchor \( \Lambda_m \) is a contraction:
\[
\|\Lambda_m\|_{\mathrm{op}} \le 1.
\]

\item[Assumption A3 (Uniformly bounded ecological memory).]
There exists a constant \( q_{\max} \ge 0 \) such that for all \( t \),
\[
\|Q_{ent}(t)\| \le q_{\max}.
\]
\end{description}

Assumption A1 is natural for bounded linear operators, and Assumption A2 holds in particular if \( \Lambda_m \) is an orthogonal projection onto a closed subspace \( \mathcal{M} \).[web:82][web:87] Assumption A3 encapsulates the effect of the finite history window and decaying kernel in the definition of \( Q_{ent} \).

\subsection{One-step norm inequality}

We first prove the basic one-step inequality used in Proposition 1 of the main text.

\begin{lemma}[One-step norm inequality]\label{lem:one-step}
Under Assumptions A1--A3, the PITN update \eqref{eq:appendix-update} satisfies
\[
\|T(t+1)\| \le c\,\|T(t)\| + q_{\max}.
\]
\end{lemma}

\begin{proof}
By Assumption A2,
\[
\|T(t+1)\|
= \|\Lambda_m(A_{p_{i_t}} T(t) + Q_{ent}(t))\|
\le \|\Lambda_m\|_{\mathrm{op}} \,\|A_{p_{i_t}} T(t) + Q_{ent}(t)\|
\le \|A_{p_{i_t}} T(t) + Q_{ent}(t)\|.
\]
Applying the triangle inequality and Assumption A1, we obtain
\[
\|T(t+1)\|
\le \|A_{p_{i_t}} T(t)\| + \|Q_{ent}(t)\|
\le \|A_{p_{i_t}}\|_{\mathrm{op}} \,\|T(t)\| + \|Q_{ent}(t)\|
\le c\,\|T(t)\| + q_{\max}.
\]
This is the desired inequality.
\end{proof}

\subsection{Affine recurrence and boundedness}

We now iterate the inequality from Lemma~\ref{lem:one-step} to derive uniform boundedness of the state norm when \( c < 1 \).

\begin{proposition}[Affine recurrence bound]\label{prop:affine-bound}
Let \( a_t = \|T(t)\| \). Under the conditions of Lemma~\ref{lem:one-step}, the sequence \( (a_t)_{t\ge 0} \) satisfies
\[
a_{t+1} \le c\,a_t + q_{\max}.
\]
If \( c < 1 \), then for all \( t \ge 0 \),
\[
a_t \le c^t a_0 + \frac{q_{\max}}{1-c} (1 - c^t),
\]
and in particular
\[
\limsup_{t\to\infty} a_t \le \frac{q_{\max}}{1-c}.
\]
\end{proposition}

\begin{proof}
The inequality \( a_{t+1} \le c\,a_t + q_{\max} \) is exactly Lemma~\ref{lem:one-step}. Solving this affine recurrence is standard: define
\[
b_t = a_t - \frac{q_{\max}}{1-c}.
\]
Then
\[
b_{t+1} = a_{t+1} - \frac{q_{\max}}{1-c}
\le c\,a_t + q_{\max} - \frac{q_{\max}}{1-c}
= c\,a_t + q_{\max} - \frac{q_{\max} - c q_{\max}}{1-c}
= c\,a_t - c\,\frac{q_{\max}}{1-c}
= c\,\left(a_t - \frac{q_{\max}}{1-c}\right)
= c\,b_t.
\]
Thus \( b_{t+1} \le c\,b_t \), implying by induction that
\[
b_t \le c^t b_0 = c^t\left(a_0 - \frac{q_{\max}}{1-c}\right).
\]
Therefore
\[
a_t 
= b_t + \frac{q_{\max}}{1-c}
\le c^t a_0 + \frac{q_{\max}}{1-c}(1 - c^t),
\]
and the claimed limit superior follows by taking \( t\to\infty \) with \( c^t \to 0 \) when \( c<1 \).
\end{proof}

As a direct corollary, we obtain the norm-stability result stated in the main body.

\begin{corollary}[Local norm stability]\label{cor:local-stability}
If \( c < 1 \), then the PITN dynamics defined by \eqref{eq:appendix-update} are norm-stable in the sense that
\[
\sup_{t\ge 0} \|T(t)\| < \infty
\]
and
\[
\|T(t)\| \to B \quad\text{with}\quad B \in \left[0,\,\frac{q_{\max}}{1-c}\right]
\]
in the sense that \( \limsup_{t\to\infty} \|T(t)\| \le \frac{q_{\max}}{1-c} \).
\end{corollary}

\begin{proof}
Immediate from Proposition~\ref{prop:affine-bound}.
\end{proof}

\subsection{Local linearization and spectral radius}

For completeness, we sketch the relation between operator-norm bounds and the spectral radius of a linearized dynamics around a fixed point, which underpins the ``spectral stability'' language in the main text.[web:87][web:121]

Let \( \bar T \in \mathcal{H} \) be a fixed point of the PITN dynamics (ignoring memory for the moment), i.e.\ it satisfies
\[
\bar T = \Lambda_m(A_{p^\ast} \bar T),
\]
for some fixed prime \( p^\ast \) or averaged operator.\footnote{In the presence of \( Q_{ent} \), one would first fix a stationary memory configuration or consider a reduced map.} Consider perturbations \( \delta T(t) = T(t) - \bar T \) and assume the update map is Fréchet differentiable at \( \bar T \). The linearization yields
\[
\delta T(t+1) \approx J\,\delta T(t),
\]
where \( J \in \mathcal{B}(\mathcal{H}) \) is the Jacobian. If we approximate
\[
J \approx \Lambda_m A_{p^\ast}
\]
and keep \( \Lambda_m \) contractive with \( \|\Lambda_m\|_{\mathrm{op}} \le 1 \), then
\[
\|J\|_{\mathrm{op}} \le \|\Lambda_m\|_{\mathrm{op}} \,\|A_{p^\ast}\|_{\mathrm{op}} \le c.
\]
Thus the spectral radius \( \rho(J) \), defined by
\[
\rho(J) := \sup\{|\lambda| : \lambda \in \sigma(J)\},
\]
satisfies \( \rho(J) \le \|J\|_{\mathrm{op}} \le c \).[web:87][web:121] Hence if \( c < 1 \), the linearized dynamics are locally asymptotically stable in the usual sense (perturbations decay in norm).

This provides a mathematically standard justification for speaking of \emph{spectral stability} or \emph{spectrally anchored} dynamics under the same operator-norm constraints used for the global norm bounds.

\section{Remarks on Non-Associativity and Path Dependence}

While composition of bounded operators in \( \mathcal{B}(\mathcal{H}) \) is associative by construction,[web:82][web:119] the PITN framework introduces non-associativity at the \emph{effective dynamics} level through two mechanisms:

\begin{enumerate}[leftmargin=*]
  \item Interleaving each application of a prime operator \( A_{p_{i_t}} \) with a projection/contraction \( \Lambda_m \) onto the variational manifold \( \mathcal{M} \), whose action need not commute with other operators.
  \item Coupling to the ecological memory term \( Q_{ent}(t) \), which depends explicitly on the history of prime words and inputs.
\end{enumerate}

Formally, for a given state \( T \), we compare
\[
T_1 = \Lambda_m\big(A_{p_2}\Lambda_m(A_{p_1} T + Q_{ent}^{(1)}) + Q_{ent}^{(2)}\big)
\]
and
\[
T_2 = \Lambda_m\big(A_{p_1}\Lambda_m(A_{p_2} T + Q_{ent}^{(1)}) + Q_{ent}^{(2)}\big),
\]
where \( Q_{ent}^{(1)}, Q_{ent}^{(2)} \) are history-dependent terms. In general, \( T_1 \neq T_2 \) even though \( A_{p_2}A_{p_1} = A_{p_2}A_{p_1} \) as formal products in \( \mathcal{B}(\mathcal{H}) \). This \emph{path dependence} is analogous to non-commuting update rules or non-ergodic dynamics in reinforcement learning, where the order of operations carries semantic information beyond the final state alone.[web:88]

\section{Summary of Bounds Used}

For rapid reference, we summarize the key operator-theoretic facts employed:

\begin{itemize}[leftmargin=*]
  \item \textbf{Operator norm.} For any \( A \in \mathcal{B}(\mathcal{H}) \), \( \|A\|_{\mathrm{op}} = \sup_{\|x\|=1} \|Ax\| \); this norm is submultiplicative: \( \|BA\|_{\mathrm{op}} \le \|B\|_{\mathrm{op}}\|A\|_{\mathrm{op}} \).[web:82][web:119]
  \item \textbf{Orthogonal projections.} Any orthogonal projection \( P \) on \( \mathcal{H} \) has \( \|P\|_{\mathrm{op}} = 1 \) and decomposes \( \mathcal{H} \) as \( \mathrm{ran}\,P \oplus \ker P \).[web:87]
  \item \textbf{Contractions.} Any operator with \( \|\Lambda_m\|_{\mathrm{op}} \le 1 \) is a contraction and does not increase vector norms.[web:120]
  \item \textbf{PITN update bound.} Under bounded prime operators and bounded memory, the PITN update obeys the affine recurrence \( \|T(t+1)\| \le c\,\|T(t)\| + q_{\max} \), yielding uniform boundedness and convergence to a ball of radius \( \frac{q_{\max}}{1-c} \) when \( c < 1 \).
  \item \textbf{Linearization and spectral radius.} For a Jacobian \( J \) satisfying \( \|J\|_{\mathrm{op}} \le c \), the spectral radius \( \rho(J) \le c \), implying local asymptotic stability when \( c<1 \).[web:87][web:121]
\end{itemize}

These results collectively justify the norm- and spectrum-based stability language in the PITN/QAGI framework while remaining within standard Hilbert-space operator theory.

\nocite{*}
\bibliographystyle{unsrt}  
\bibliography{references}
\end{document}